% WHAT \showlists SAYS OF A FORMULA.
%
% A formula is a list of its own and a level of the nest of its own, and
% \showlists is the only thing that ever writes it: a formula becomes a
% horizontal list at its closing shift, so \showbox never sees one.
%
%   ### math mode entered at line 10
%   \mathord
%   .\fam1 A
%   \mathord
%   .\fam0 +
%
% Every noad stands on a line of its own. The three fields an atom may have
% are written under it with a marker of their own -- `.' for the nucleus,
% `^' for the superscript, `_' for the subscript -- and that marker is also
% the recursion history \showboxdepth counts:
%
%   \mathord         an atom, named for its class
%   \vcenter         an atom \vcenter made, named for what it is
%   \mathop\limits   the limit controls, where they are not the default
%   \overline        \underline
%   \radical"161     the delimiter, as its twenty-four bit code in hex
%   \accent\fam0 ^^V a \mathaccent, named for the character it puts on
%   \displaystyle    a style, and its three fellows
%   \mkern3.0mu      \glue(\mskip) 4.0mu plus 1.0fil     \glue(\nonscript)
%   \mathchoice      with its four lists under D, T, S and s
%   \left"0 ...      the delimiter a \left group opened, at the head of it
%
% A \fam and a character make a field: `.\fam1 A'. A field that is a
% sub-formula is written as the LIST it is, not as any box that was made of
% it, and one that came out empty is `{}'. A glue, a kern, a penalty or a
% box written in a formula is a node like any other and is written as one.
%
% A sub-formula that stands as an ORDINARY ATOM is the nucleus of a noad
% that is already in the list around it, so that list ends with a bare
% `\mathord' -- or `\mathop', or whatever class was forced on it -- while
% the braces are open. A sub-formula filling a script, a radicand, a
% \mathchoice branch or a \vcenter fills a noad that is there already, or
% none at all, and adds nothing.
%
% An \over makes the whole list a fraction: what stands after it is the
% denominator so far, and the fraction waiting for the rest is written under
% `this will begin denominator of:' with its numerator behind a `\' and its
% denominator behind a `/'. A numerator is a list from the moment the \over
% is read, so an empty one is `{}'; a denominator still waiting for its list
% writes nothing.
%
% HSTeX wrote no formula at all: a formula was not a level of the nest, and
% the level under it was stamped with the formula's own mode instead. trip
% asks at lines 284 and 439.

\catcode`\{=1 \catcode`\}=2 \catcode`\$=3 \catcode`\^=7 \catcode`\_=8
\tracingonline=1 \showboxbreadth=99 \showboxdepth=99
\font\rip=trip \fontdimen22\rip=7pt \rip
\textfont0=\rip \scriptfont0=\rip \scriptscriptfont0=\rip
\textfont1=\rip \scriptfont1=\rip \scriptscriptfont1=\rip
\textfont2=\rip \scriptfont2=\rip \scriptscriptfont2=\rip
\textfont3=\rip \scriptfont3=\rip \scriptscriptfont3=\rip
\message{[A a plain formula]}
\setbox0=\hbox{$A+B\showlists$}
\message{[B a braced sub-formula]}
\setbox0=\hbox{$A{B\showlists}$}
\message{[C a script]}
\setbox0=\hbox{$A^{B\showlists}$}
\message{[D a fraction waiting for its denominator]}
\setbox0=\hbox{$A\over B\showlists$}
\message{[E classes, styles and mu glue]}
\setbox0=\hbox{$\mathbin{A}\scriptstyle\mkern3mu\mskip4mu plus1fil\nonscript B\showlists$}
\message{[F an overline, an underline and a radical]}
\setbox0=\hbox{$\overline A\underline B\radical"161 C\showlists$}
\message{[G a mathchoice]}
\setbox0=\hbox{$\mathchoice{A}{B}{C}{D}\showlists$}
\message{[H an accent and a vcenter]}
\setbox0=\hbox{$\mathaccent"7016 A\vcenter{\hbox{}}\showlists$}
\message{[I a left group]}
\setbox0=\hbox{$\left(A\showlists\right)$}
\message{[J a box and a node in a formula]}
\setbox0=\hbox{$A\hbox{x}\penalty5 \kern3pt\showlists$}
\message{[K what a group of each kind leaves in the list around it]}
\setbox0=\hbox{$X{A\showlists}$}
\setbox0=\hbox{$X\mathop{A\showlists}$}
\setbox0=\hbox{$X\mathbin{A\showlists}$}
\setbox0=\hbox{$X\overline{A\showlists}$}
\setbox0=\hbox{$X\underline{A\showlists}$}
\setbox0=\hbox{$X\radical"161{A\showlists}$}
\setbox0=\hbox{$X\mathaccent"7016{A\showlists}$}
\setbox0=\hbox{$X\vcenter{\hbox{}\showlists}$}
\setbox0=\hbox{$X^{B}_{C\showlists}$}
\setbox0=\hbox{$X\mathchoice{A\showlists}{}{}{}$}
\setbox0=\hbox{$X{Y{A\showlists}}$}
\message{[done]}
\end
